\section{2026 年 11 月核验增补}

本节按论文首次提交时间排列。每项方法附有可核验的外部入选依据。

% METHOD: 2510.03312
\paragraph{Universal Beta Splatting：用形状可调的高维核统一空间与时角变化}
2025.10~\cite{2510_03312}。UBS 用各向异性的高维 Beta 核替代固定的高斯轮廓。基础三维位置之外，额外维度可表示时间或观测方向；核形状参数调节这些方向的响应，空间正交参数化和条件切片将当前查询转为可渲染原语。在相应参数极限下，它可以近似高斯核。

原文在 NeRF Synthetic、Mip-NeRF 360、6DGS-PBR、D-NeRF 和 7DGS-PBR 评价静态、动态及强角度依赖。静态与动态实验使用不同硬件，论文也注明部分工作比较最后一次迭代与原报告的最佳迭代。体渲染场景的最大提升不能当作所有场景平均。作者页列有 United Imaging Intelligence。

\textbf{入选依据} 论文署名单位包括 United Imaging Intelligence（\href{https://arxiv.org/html/2510.03312v3}{论文原文}）。

% METHOD: 2607.26578
\paragraph{3DGBGS：用点云自适应粒球确定锚点与尺度}
2026.07~\cite{2607_26578}。3DGBGS 用极粗或极细的一套固定体素划分都不理想，改为按 SfM 点云的局部结构划分多粒度球。球心初始化锚点，半径提供局部尺度先验，粗尺度覆盖较平滑区域，小粒球保留更复杂几何。

论文在十九个真实新视角合成场景比较，在维持可比图像质量时降低锚点和模型存储。表示收益基于给定 SfM 输入，不直接说明未标定或稀疏位姿情况下也能保真。作者 Yi Wang 在论文中列出的单位是重庆蚂蚁消费金融有限公司。

\textbf{入选依据} 论文署名单位包括 重庆蚂蚁消费金融有限公司（\href{https://arxiv.org/html/2607.26578v1}{论文原文}）。
